% Options for packages loaded elsewhere
\PassOptionsToPackage{unicode}{hyperref}
\PassOptionsToPackage{hyphens}{url}
\PassOptionsToPackage{dvipsnames,svgnames,x11names}{xcolor}
\documentclass[
  10pt,
]{article}
\usepackage{xcolor}
\usepackage[margin=2.2cm]{geometry}
\usepackage{amsmath,amssymb}
\setcounter{secnumdepth}{-\maxdimen} % remove section numbering
\usepackage{iftex}
\ifPDFTeX
  \usepackage[T1]{fontenc}
  \usepackage[utf8]{inputenc}
  \usepackage{textcomp} % provide euro and other symbols
\else % if luatex or xetex
  \usepackage{unicode-math} % this also loads fontspec
  \defaultfontfeatures{Scale=MatchLowercase}
  \defaultfontfeatures[\rmfamily]{Ligatures=TeX,Scale=1}
\fi
\usepackage{lmodern}
\ifPDFTeX\else
  % xetex/luatex font selection
\fi
% Use upquote if available, for straight quotes in verbatim environments
\IfFileExists{upquote.sty}{\usepackage{upquote}}{}
\IfFileExists{microtype.sty}{% use microtype if available
  \usepackage[]{microtype}
  \UseMicrotypeSet[protrusion]{basicmath} % disable protrusion for tt fonts
}{}
\makeatletter
\@ifundefined{KOMAClassName}{% if non-KOMA class
  \IfFileExists{parskip.sty}{%
    \usepackage{parskip}
  }{% else
    \setlength{\parindent}{0pt}
    \setlength{\parskip}{6pt plus 2pt minus 1pt}}
}{% if KOMA class
  \KOMAoptions{parskip=half}}
\makeatother
\usepackage{graphicx}
\makeatletter
\newsavebox\pandoc@box
\newcommand*\pandocbounded[1]{% scales image to fit in text height/width
  \sbox\pandoc@box{#1}%
  \Gscale@div\@tempa{\textheight}{\dimexpr\ht\pandoc@box+\dp\pandoc@box\relax}%
  \Gscale@div\@tempb{\linewidth}{\wd\pandoc@box}%
  \ifdim\@tempb\p@<\@tempa\p@\let\@tempa\@tempb\fi% select the smaller of both
  \ifdim\@tempa\p@<\p@\scalebox{\@tempa}{\usebox\pandoc@box}%
  \else\usebox{\pandoc@box}%
  \fi%
}
% Set default figure placement to htbp
\def\fps@figure{htbp}
\makeatother
\setlength{\emergencystretch}{3em} % prevent overfull lines
\providecommand{\tightlist}{%
  \setlength{\itemsep}{0pt}\setlength{\parskip}{0pt}}
\usepackage{fontspec}
\setmainfont{STIX Two Text}
\setmonofont{Menlo}[Scale=0.85]
\usepackage{xeCJK}
\setCJKmainfont{Songti SC}
\usepackage{microtype}
\usepackage{etoolbox}
\AtBeginEnvironment{longtable}{\footnotesize}
\usepackage{xurl}
\setlength{\emergencystretch}{3em}
\sloppy
\setCJKmonofont{Songti SC}
\usepackage{unicode-math}
\setmathfont{STIX Two Math}
\usepackage{newunicodechar}
\newunicodechar{→}{\ensuremath{\rightarrow}}
\newunicodechar{∎}{\ensuremath{\blacksquare}}
\newunicodechar{≠}{\ensuremath{\neq}}
\newunicodechar{≈}{\ensuremath{\approx}}
\newunicodechar{↔}{\ensuremath{\leftrightarrow}}
\newunicodechar{≪}{\ensuremath{\ll}}
\newunicodechar{↛}{\ensuremath{\nrightarrow}}
\newunicodechar{・}{\textperiodcentered}
\newunicodechar{⇒}{\ensuremath{\Rightarrow}}
\newunicodechar{∪}{\ensuremath{\cup}}
\newunicodechar{≥}{\ensuremath{\geq}}
\newunicodechar{≤}{\ensuremath{\leq}}
\newunicodechar{⟨}{\ensuremath{\langle}}
\newunicodechar{⟩}{\ensuremath{\rangle}}
\usepackage{bookmark}
\IfFileExists{xurl.sty}{\usepackage{xurl}}{} % add URL line breaks if available
\urlstyle{same}
% fallback for those not using the hyperref driver hyperxmp:
\makeatletter
\@ifundefined{xmpquote}{\newcommand{\xmpquote}[1]{#1}}{}
\makeatother
\hypersetup{
  pdftitle={The Price of the Unexamined Interface},
  pdfauthor={David Charlot, Open Interface Engineering},
  colorlinks=true,
  linkcolor={blue},
  filecolor={Maroon},
  citecolor={Blue},
  urlcolor={blue},
  pdfcreator={LaTeX via pandoc}}

\title{The Price of the Unexamined Interface}
\author{David Charlot, Open Interface Engineering}
\date{Version 1.1, 11 October 2026. HTML:
https://research.openie.dev/papers/interface-economics/}

\begin{document}
\maketitle

{
\setcounter{tocdepth}{2}
\tableofcontents
}
\emph{Version 1.0, 11 October 2026. LaTeX PDF:
\href{https://research.openie.dev/pdfs/latex/interface-economics.pdf}{/pdfs/latex/interface-economics.pdf}
(\href{https://research.openie.dev/pdfs/latex/interface-economics.tex}{source}).}

\emph{The economics companion to
\href{https://research.openie.dev/papers/three-body/}{The Pursuit of AGI
Resolves the Three-Body Problem in Computing}. The research paper holds
the engineering and the theory. This essay is where the interface thesis
meets prices, budgets and markets.}

\subsection{Abstract}\label{abstract}

The research paper argues that the instability of computing systems
lives in the interfaces between three bodies, theorists, makers and
users, where assumptions cross unwritten and unchecked. This essay asks
what that instability costs, and answers only with figures that a
primary source states. Each figure is labeled by kind. An
\textbf{accounted cost} is what an agency or firm reports it actually
spent or lost. A \textbf{model estimate} is a figure produced by a
stated method from surveys, interviews or secondary data. A
\textbf{projection} is a statement about the future. The three kinds are
not added together.

The accounted costs of single interface failures are large and specific:
the Mars Surveyor 1998 project, which included Mars Climate Orbiter,
cost \$327.6 million {[}1{]}, and ESA approved 214 million ECU to
recover the Cluster mission lost on the first Ariane 5 {[}4{]}. The
model estimates are larger and broader: NIST put the annual cost of an
inadequate software testing infrastructure in the United States at
\$59.5 billion {[}6{]}, and the cost of inadequate interoperability in
the US capital facilities industry at \$15.8 billion a year {[}7{]}.
CISQ estimated the cost of poor software quality in the United States at
no less than \$2.41 trillion in 2022 and accumulated technical debt at
about \$1.52 trillion {[}8{]}. None of these measures interface failure
alone. Each contains it.

\subsection{1. Kinds of number}\label{kinds-of-number}

Economic claims about software failure travel in three forms that look
alike and are not. A project budget is an accounting record. A national
cost estimate is the output of a model whose inputs are surveys, case
studies and assumptions about scale. A forecast is neither. The
\href{https://research.openie.dev/papers/ai-coding-market/\#1-projections-and-measurements-are-different-objects}{AI-coding-market
essay} makes the same separation for the AI coding market, and this
essay follows its rule: every figure carries its kind, its source and
what it actually counts.

Open Interface Engineering's public site states that more than \$1
trillion is lost annually to interface and integration failures
{[}17{]}. That sentence is a headline, not a primary estimate, and the
site gives no method for it. The nearest primary estimates are listed in
Section 3. None of them isolates interfaces, so this essay does not map
them onto that figure or onto any single total.

\subsection{2. Single failures,
accounted}\label{single-failures-accounted}

\textbf{Mars Climate Orbiter.} The interface was a unit: thruster
impulse data in pound-force seconds read as newton seconds {[}2{]}.
NASA's project fact sheet gives the cost of the Mars Surveyor 1998
project as \$327.6 million for the orbiter and the Mars Polar Lander
together: \$193.1 million for spacecraft development, \$91.7 million for
launch and \$42.8 million for operations {[}1{]}. Contemporary reports
gave the orbiter alone as a \$125 million spacecraft {[}3{]}. Kind:
accounted cost of the project, with a reported figure for the one
spacecraft.

\textbf{Ariane 501 and Cluster.} The interface was an envelope: software
reused from Ariane 4 assumed a range of horizontal velocity that the new
launcher exceeded {[}5{]}. The four Cluster satellites were lost. ESA's
account of the recovery records the options its Science Programme
Committee priced in 1996, including 351 million accounting units to
rebuild the four satellites to the original design, and ESA approved a
recovery mission, Cluster II, for 214 million ECU in 1997 {[}4, 18{]}.
Kind: accounted cost of the recovery that ESA chose, which is a lower
bound on the loss, not the value of the satellites or of the delayed
science.

These two cases fix the shape of the problem. In each, every body did
its own work, and the cost landed at the interface. In each, the check
that would have caught the mismatch was cheap relative to the loss.
Neither case can be scaled up to a national figure, because single
accidents are not a sample.

\subsection{3. National estimates,
modeled}\label{national-estimates-modeled}

\includegraphics[width=0.95\linewidth,height=\textheight,keepaspectratio]{fig/estimates-by-kind.png}

\emph{Figure 1. The estimates by kind. Data figure, logarithmic axis.
Values as printed by the sources {[}1, 4, 6, 7, 8, 9{]}. The bars are
not additive: the CISQ total contains categories that overlap the NIST
and Stripe figures, and a stock (technical debt) is drawn beside annual
flows. A
\href{https://research.openie.dev/figures/interface-economics/estimates-by-kind.png}{PNG
copy} is provided.}

\textbf{Software testing, NIST 2002.} RTI International, for NIST,
estimated that an inadequate infrastructure for software testing cost
the US economy \$22.2 billion to \$59.5 billion a year, with the upper
figure as the total cost and the lower figure as the part that feasible
improvements could remove {[}6{]}. The method combined surveys of
developers and users in two sectors, transportation manufacturing and
financial services, with an extrapolation to the national economy. Kind:
model estimate.

\textbf{Interoperability, NIST 2004.} Gallaher and colleagues, for NIST,
estimated the cost of inadequate interoperability in the US capital
facilities industry at \$15.8 billion a year, from interviews and
surveys of owners, operators, architects, engineers and contractors
{[}7{]}. Interoperability here is the ability to exchange and use design
and operating data across software systems, which is an interface in
exactly the paper's sense. The report states that its figure omits costs
it could not quantify. Kind: model estimate, one industry, presented as
conservative.

\textbf{Poor software quality, CISQ 2022.} The Consortium for
Information and Software Quality estimated that the cost of poor
software quality in the United States had grown to at least \$2.41
trillion in 2022, and that accumulated technical debt had grown to about
\$1.52 trillion {[}8{]}. The report builds the total from categories,
operational failures, unsuccessful projects, legacy system problems and
cybercrime losses attributed to software vulnerabilities, and draws much
of its input from secondary sources; its cybercrime component apportions
a global forecast to the United States {[}8{]}. Kind: model estimate
built partly on other estimates. It is the largest documented figure in
this area, and it is a figure for software quality as a whole, not for
interfaces.

\textbf{Developer time, Stripe 2018.} A Stripe survey conducted with
Harris Poll reported that developers spend about 17.3 hours a week on
maintenance work such as debugging and refactoring, and about four hours
a week on ``bad code,'' which the report values at nearly \$85 billion a
year worldwide {[}9{]}. Kind: survey-based estimate.

What these figures share is a structure, not a total. Each counts the
cost of an assumption that was relied on and not checked: an untested
path, a data format that does not carry across, a vulnerability at a
boundary, a legacy system whose contracts no one can restate. The
research paper's
\href{https://research.openie.dev/papers/three-body/\#52-propositions}{Proposition
2} explains why such costs show up far from their cause, which is also
why national estimates must be modeled rather than counted.

\subsection{4. Technical debt is an interface
account}\label{technical-debt-is-an-interface-account}

Cunningham introduced technical debt in 1992 as the cost of shipping
code whose structure does not yet match what the team has learned, with
interest paid on every later change until the debt is repaid {[}10{]}.
Sculley and colleagues showed that in machine learning systems most of
that debt sits outside the model, in glue code, configuration, data
dependencies and undeclared consumers {[}11{]}. Read through the
research paper, technical debt is the balance of unexamined interface
assumptions carried on the books. Its interest is Proposition 2: an
unchecked mismatch charges every level above it, on every run.

CISQ's estimate of about \$1.52 trillion of accumulated technical debt
in the United States is a stock, measured in the cost to fix {[}8{]}.
The annual figures in Section 3 are flows. The research paper's
Proposition 5 states the economic consequence in the units of the
companion papers: a check paid once at a boundary costs less over a
workload than a failure paid again at every level, and the difference
compounds with use.

\subsection{5. Standards are how markets close
interfaces}\label{standards-are-how-markets-close-interfaces}

Economists have studied the interface problem for decades under another
name: compatibility. Katz and Shapiro showed that when a product's value
rises with the number of users it can work with, firms' incentives to be
compatible depend on their size and position, and the market can deliver
too little compatibility {[}12{]}. Farrell and Saloner showed that a
market can stay stuck on an old standard, or move to a new one too fast,
because each user's choice depends on what others choose {[}13{]}.
David's history of the QWERTY keyboard made path dependence concrete
{[}14{]}. Baldwin and Clark showed that design rules, visible interfaces
fixed in advance, are what let modular industries form at all {[}15{]}.
Swann's review for the UK government summarized evidence that standards
contribute to productivity and trade {[}16{]}.

The common thread is that an interface is a public good with private
costs. Writing an assumption down and checking it benefits every party
that crosses the boundary, but it is paid for by whoever does the
writing. That is the economic form of the research paper's observation
that each body assumes the handoff is someone else's job. Standards
bodies, regulators and dominant platforms are the institutions that have
historically paid that cost. A representation in which every interface
assumption is explicit and checked by construction, as in the
\href{https://research.openie.dev/papers/rce/}{Reality Code Engine},
lowers the private cost of closing an interface toward the cost of
running its check.

\subsection{6. Ties to the AI coding
market}\label{ties-to-the-ai-coding-market}

The
\href{https://research.openie.dev/papers/ai-coding-market/}{AI-coding-market
essay} sets out the projections for AI coding: about \$3 trillion a year
in developer value, company revenue, token volume and capital spending,
each labeled by kind. Those projections price the generation of code.
They do not price its interfaces. Every generated change crosses the
same boundaries as hand-written code, between a specification and an
implementation, between a library and its users, between software and
the instruction contract beneath it, and each crossing carries
assumptions that the generator does not check.

The figures in Section 3 suggest where the value of verification sits.
If the cost of poor software quality is of the order CISQ estimates, and
a large share of it is the cost of unchecked assumptions at boundaries,
then a method that closes interfaces is worth more than a method that
writes code faster. That is a proposition about where value lies, not a
measurement. The protocols of the research paper
(\href{https://research.openie.dev/papers/three-body/\#8-protocols}{Section
8}) and of the AI-coding-market essay are the way to measure it.

\subsection{7. Propositions and
protocols}\label{propositions-and-protocols}

\textbf{Proposition E1 (the interface share).} In a population of
software failures with accounted costs, the share of cost whose cause is
a violated interface assumption, in the sense of the research paper's
Definition 3, exceeds the share whose cause lies wholly inside one body.

\emph{Protocol.} Take a public incident corpus with cost fields, such as
agency mishap reports or regulatory filings. Two coders classify each
cause as interface or internal using the paper's definitions, blind to
cost. Report the cost-weighted interface share with a bootstrap interval
and the coders' agreement.

\textbf{Proposition E2 (checks pay).} For a fixed workload, the cost of
adding an explicit check at an interface is lower than the expected cost
of the failures it refuses, summed over every level the failure would
otherwise reach.

\emph{Protocol.} Run the research paper's break and catch protocol
(\href{https://research.openie.dev/papers/three-body/\#82-break-and-catch}{§8.2})
on a system with priced levels, recording the cost of each check and the
cost of each uncaught violation where it lands. Report both in joules
and in dollars at a stated price.

\textbf{Proposition E3 (closing an interface lowers its private cost).}
When the interface assumption is part of a shared representation and
checked by construction, the cost to a single party of closing it falls
toward the cost of running the check.

\emph{Protocol.} The cross-body handoff protocol of the research paper
(\href{https://research.openie.dev/papers/three-body/\#84-cross-body-handoff}{§8.4}),
with labor hours recorded on each side of the handoff in both arms.

\subsection{8. Conclusion}\label{conclusion}

The cost of unexamined interfaces is real, large and badly measured.
Single failures can be priced exactly, and they show that the loss lands
at a boundary every body thought belonged to someone else. National
estimates are models, and the largest of them describe software quality
as a whole. No primary source states the cost of interface failure by
itself, and this essay does not supply one. What the evidence supports
is a structure: costs that arise at boundaries, propagate upward and
compound as debt, in markets that underprovide the public good of a
checked interface. The research paper gives that structure its theory
and its tests.

\subsection{References}\label{references}

{[}1{]} NASA Jet Propulsion Laboratory, ``Mars Climate Orbiter Fact
Sheet,'' Mars Surveyor 1998 project. Archived copy:
https://web.archive.org/web/20121003073752/http://mars.jpl.nasa.gov/msp98/orbiter/fact.html

{[}2{]} Mars Climate Orbiter Mishap Investigation Board, \emph{Phase I
Report}. NASA, 10 November 1999.
https://llis.nasa.gov/llis\_lib/pdf/1009464main1\_0641-mr.pdf

{[}3{]} A. Pollack, ``Two Teams, Two Measures Equaled One Lost
Spacecraft,'' \emph{The New York Times}, 1 October 1999.
https://archive.nytimes.com/www.nytimes.com/library/national/science/100199sci-nasa-mars.html

{[}4{]} European Space Agency, ``Cluster will fly again!,'' press
release N° 12 1997, 26 March 1997.
https://www.esa.int/Newsroom/Press\_Releases/Cluster\_will\_fly\_again

{[}5{]} J. L. Lions (chair), \emph{Ariane 5 Flight 501 Failure: Report
by the Inquiry Board}. European Space Agency, 19 July 1996.
https://esamultimedia.esa.int/docs/esa-x-1819eng.pdf

{[}6{]} RTI International for the National Institute of Standards and
Technology, \emph{The Economic Impacts of Inadequate Infrastructure for
Software Testing}, Planning Report 02-3, May 2002.
https://www.nist.gov/system/files/documents/director/planning/report02-3.pdf

{[}7{]} M. P. Gallaher, A. C. O'Connor, J. L. Dettbarn and L. T. Gilday,
\emph{Cost Analysis of Inadequate Interoperability in the U.S. Capital
Facilities Industry}, NIST GCR 04-867, 2004.
\href{https://doi.org/10.6028/NIST.GCR.04-867}{doi:10.6028/NIST.GCR.04-867}

{[}8{]} H. Krasner, \emph{The Cost of Poor Software Quality in the US: A
2022 Report}. Consortium for Information and Software Quality, 2022.
https://www.it-cisq.org/the-cost-of-poor-quality-software-in-the-us-a-2022-report/

{[}9{]} Stripe, \emph{The Developer Coefficient}, with Harris Poll,
September 2018.
https://stripe.com/files/reports/the-developer-coefficient.pdf

{[}10{]} W. Cunningham, ``The WyCash portfolio management system,''
\emph{OOPSLA 1992 Addendum}, \emph{ACM SIGPLAN OOPS Messenger} 4(2),
1993, pp.~29 to 30.
\href{https://doi.org/10.1145/157710.157715}{doi:10.1145/157710.157715}

{[}11{]} D. Sculley, G. Holt, D. Golovin, E. Davydov, T. Phillips, D.
Ebner, V. Chaudhary, M. Young, J.-F. Crespo and D. Dennison, ``Hidden
Technical Debt in Machine Learning Systems,'' \emph{Advances in Neural
Information Processing Systems} 28, 2015.
https://papers.nips.cc/paper\_files/paper/2015/hash/86df7dcfd896fcaf2674f757a2463eba-Abstract.html

{[}12{]} M. L. Katz and C. Shapiro, ``Network Externalities,
Competition, and Compatibility,'' \emph{American Economic Review} 75(3),
1985, pp.~424 to 440. https://www.jstor.org/stable/1814809

{[}13{]} J. Farrell and G. Saloner, ``Standardization, Compatibility,
and Innovation,'' \emph{RAND Journal of Economics} 16(1), 1985, pp.~70
to 83. \href{https://doi.org/10.2307/2555589}{doi:10.2307/2555589}

{[}14{]} P. A. David, ``Clio and the Economics of QWERTY,''
\emph{American Economic Review} 75(2), 1985, pp.~332 to 337.
https://www.jstor.org/stable/1805621

{[}15{]} C. Y. Baldwin and K. B. Clark, \emph{Design Rules, Volume 1:
The Power of Modularity}. MIT Press, 2000.
\href{https://doi.org/10.7551/mitpress/2366.001.0001}{doi:10.7551/mitpress/2366.001.0001}

{[}16{]} G. M. P. Swann, \emph{The Economics of Standardization: An
Update}. Report for the UK Department for Business, Innovation and
Skills, 2010.
https://www.gov.uk/government/publications/economics-of-standardisation-update-to-report

{[}17{]} Open Interface Engineering, homepage, accessed 11 October 2026.
https://openie.dev/

{[}18{]} European Space Agency, ``The Resurrection of the Cluster
Scientific Mission,'' \emph{ESA Bulletin} 91, 1997.
https://www.esa.int/esapub/bulletin/bullet91/b91cred.htm

\end{document}
